\magicSection{Feature Wishlist}{sec:Features}

As mentioned, the initial project will focus on adapting Exo-GPU ideas into a new form more capable of expressing maximum performance kernels.
Primarily, this requires adding \myKeyA{mutable pointer} variables, and replacing \emph{checked} rewrites with a new functional checking system.
The new system needs to be capable of reasoning about mutable pointers, and fine-grained reordering and ring buffering that I was unable to express using Exo rewrites.
We will retain the ``abstract machine'' synchronization checking from Exo-GPU.

We will still be gambling on \myKeyA{internal determinism}, or ``sequential-parallel equivalence''.
This factors correctness checking into two problems: \myKeyA{value correctness}, which is whether the algorithm, with all parallel loops and async instructions serialized, will implement the functional specification, and \myKeyA{sync correctness}, which is whether the real parallel algorithm is sufficiently synchronized to be functionally equivalent to the serialized version of the algorithm.

This sets up a chain of equivalence

\begin{center}
    Functional spec $\xrightarrow{\text{\myKeyA{value correctness}}}$
    Serialized algorithm $\xrightarrow{\text{\myKeyA{sync correctness}}}$
    Parallelized algorithm
\end{center}

Exo-GPU supports one ``escape hatch'' from determinism, which is re-ordering of non-associative floating point atomics.
In principle, we could support more escape hatches, e.g.
\begin{itemize}
  \item Mutex-guarded regions, with commutative operations within (or any operations if we don't care about determinism at all).
  \item Opportunistic data-reuse: reading a value that may or may not be ready (e.g. dynamic programming, KV cache with eviction, decoupled look-back), which only requires analyzing 2 cases.
  \item Some opaque spec function \lighttt{reorder(i)} that is an unknown bijection from integers to integers, representing the unknown order that \lighttt{atomicAdd(\&..., 1)} will give unique values to threads.
\end{itemize}

Anyway, the features I have in mind are

\begin{itemize}
  \item Retaining Exo-GPU's model of sequential and parallel loops and CUDA kernel launch \textbf{(Section~\ref{sec:SeqPar})}
  \begin{itemize}
    \item Static mapping of parallel ``loop iterations'' to threads
    \item Static typing of collective units (thread groupings) that execute statements (e.g. thread, warp cooperative, CTA cooperative)
    \item See \lighttt{spork\_b.pdf} collective tiling chapter
  \end{itemize}
  \item Tensor variables with ``typed strides'' \textbf{(Section~\ref{sec:Tensors})}
  \begin{itemize}
    \item Each dimension is annotated with a stride that is in units of threads, CTAs, RAM address, register ID, or bits (for multiple values packed in one word)
    \item Extends Exo-GPU's binary distinction between distributed dimensions (strided by threads) and non-distributed dimensions (strided by anything else)
    \item Like in Exo-GPU, thread shards/fragments are explicit, not implicitly duplicated like in CUDA C++.
  \end{itemize}
  \item Window references (mutable ``pointer'' variables) \textbf{(Section~\ref{sec:Windows})}
  \item Polynomial-based functional specifications \textbf{(Section~\ref{sec:PolynomialSpec})}
  \item Annotations on assignment statements that specify, as a function of control variables and the functional specification \textbf{(Section~\ref{sec:AssignmentAnnotation})}
  \begin{itemize}
    \item the tensor elements referenced by any window references used \textbf{(Section~\ref{sec:WindowAnnotation})}
    \item what values are read \textbf{(Section~\ref{sec:ReadAnnotation})}
    \item what value is written \textbf{(Section~\ref{sec:WriteAnnotation})}
  \end{itemize}
  This is just a seed idea to start on pragmatic functional correctness checking; as I mentioned, I don't know much about formal methods
  \item Value correctness checking for concrete problem sizes \textbf{(Section~\ref{sec:ValueCorrectnessChecking})}
  \item Sync correctness checking like in Exo-GPU \textbf{(Section~\ref{sec:SyncCheck})}
  \item Explicit instruction and procedure calls \textbf{(Section~\ref{sec:Proc})}
\end{itemize}
